\documentclass[10pt,aspectratio=169]{beamer}
\usetheme{Madrid}
\usecolortheme{whale}

\usepackage[utf8]{inputenc}
\usepackage[french]{babel}
\usepackage{graphicx}
\usepackage{booktabs}
\usepackage{xcolor}
\usepackage{amsmath}
\usepackage{tikz}
\usetikzlibrary{shapes,arrows,positioning}

\title[Architecture \& Workflow VC-UY1]{Architecture Globale et Workflow d'Apprentissage Hybride}
\subtitle{Intégration de RLS-ARX, GRU-EWC et Quantification TinyML pour le Calcul Volontaire (VC-UY1)}
\author[Mémoire de Master 2]{Présenté par : \textbf{Neussi} \\[0.5em] \small Encadrant : \textbf{Dr. Adamou Hamza}}
\institute[UY1]{Université de Yaoundé I \\ Département d'Informatique \\ Spécialité : Systèmes et Réseaux}
\date[Juin 2026]{Soutenance de Travaux de Recherche}

\begin{document}

% Slide de Titre
\begin{frame}
  \titlepage
\end{frame}

% Slide Plan
\begin{frame}{Plan de la Présentation}
  \begin{columns}[T]
    \begin{column}{0.48\textwidth}
      \begin{block}{1. Architecture Système}
        \begin{itemize}
          \item Les 3 Piliers du système VC-UY1
          \item Structure de l'agent local sur le nœud
        \end{itemize}
      \end{block}
      
      \begin{block}{2. Le Modèle Hybride}
        \begin{itemize}
          \item Branche RLS-ARX (Linéaire)
          \item Branche GRU-EWC (Non-linéaire)
        \end{itemize}
      \end{block}
    \end{column}
    
    \begin{column}{0.48\textwidth}
      \begin{block}{3. Workflow d'Entraînement}
        \begin{itemize}
          \item Pré-entraînement global (SETI@home)
          \item Quantification TinyML (INT8)
          \item Co-adaptation et régularisation locale
        \end{itemize}
      \end{block}
      
      \begin{block}{4. Intégration \& Ordonnancement}
        \begin{itemize}
          \item Soumission et migration de tâches
        \end{itemize}
      \end{block}
    \end{column}
  \end{columns}
\end{frame}

% Slide 3: Les 3 Piliers de l'Architecture
\begin{frame}{1. L'Architecture Générale en 3 Piliers}
  \begin{columns}[T]
    \begin{column}{0.58\textwidth}
      \begin{block}{Description des Composants}
        \begin{itemize}
          \item \textbf{Pilier 1 : Télémétrie Frugale} \\
                Collecte locale (60s) de l'alimentation, batterie, réseau et utilisateur. Buffer glissant de 72h.
          \item \textbf{Pilier 2 : Baseline Linéaire RLS-ARX} \\
                Modèle adaptatif léger à chaque minute ($< 1$ Ko RAM) pour réagir aux délestages.
          \item \textbf{Pilier 3 : Modèle Profond GRU-EWC} \\
                Pré-entraînement global puis adaptation locale anti-oubli.
        \end{itemize}
      \end{block}
    \end{column}
    
    \begin{column}{0.38\textwidth}
      \begin{exampleblock}{Schéma de l'Architecture}
        \centering
        \vspace{0.5em}
        \begin{tikzpicture}[scale=0.7, every node/.style={transform shape},
            box/.style={rectangle, draw=blue!50!black, fill=blue!10, rounded corners, minimum width=3cm, minimum height=1cm, align=center, font=\scriptsize, text width=2.8cm},
            sched/.style={rectangle, draw=green!50!black, fill=green!10, rounded corners, minimum width=3.3cm, minimum height=0.8cm, align=center, font=\scriptsize\bfseries}]
            
            \node[box] (p1) {Pilier 1: Télémétrie\\{\tiny Collecte \& Buffer 72h}};
            \node[box] (p2) [below=0.6cm of p1] {Pilier 2: RLS-ARX\\{\tiny Baseline Linéaire}};
            \node[box] (p3) [below=0.6cm of p2] {Pilier 3: GRU-EWC\\{\tiny Modèle Profond}};
            \node[sched] (sch) [right=0.8cm of p2] {Ordonnanceur\\{\tiny Tasks VC-UY1}};
            
            \draw[-latex, thick] (p1) -- (p2);
            \draw[-latex, thick] (p2) -- (p3);
            \draw[-latex, thick] (p2) -- (sch);
            \draw[-latex, thick] (p3) -| (sch);
        \end{tikzpicture}
      \end{exampleblock}
    \end{column}
  \end{columns}
\end{frame}

% Slide 4: L'Architecture du Prédicteur Hybride
\begin{frame}{2. L'Architecture Interne de l'Agent Prédictif}
  \begin{columns}[T]
    \begin{column}{0.52\textwidth}
      \begin{block}{Co-existence des deux branches}
        \begin{itemize}
          \item Les deux modèles reçoivent en parallèle les mesures d'état courantes.
          \item \textbf{Branche RLS-ARX :} Agit comme un réflexe local rapide pour détecter les chutes de courant immédiates.
          \item \textbf{Branche GRU-EWC :} Agit comme le cerveau prédisant les tendances quotidiennes de l'utilisateur.
        \end{itemize}
      \end{block}
    \end{column}
    
    \begin{column}{0.44\textwidth}
      \begin{exampleblock}{Flux de Prédiction Hybride}
        \centering
        \vspace{0.2em}
        \begin{tikzpicture}[scale=0.65, every node/.style={transform shape},
            input/.style={circle, draw=blue!50!black, fill=blue!10, minimum size=0.8cm, font=\scriptsize},
            branch/.style={rectangle, draw=blue!50!black, fill=blue!5, rounded corners, minimum width=3.5cm, minimum height=0.8cm, align=center, font=\scriptsize},
            merge/.style={circle, draw=orange!50!black, fill=orange!10, minimum size=0.8cm, font=\scriptsize\bfseries},
            output/.style={rectangle, draw=green!50!black, fill=green!10, rounded corners, minimum width=2.5cm, minimum height=0.6cm, align=center, font=\scriptsize}]

            \node[input] (in) {Input $x_t$};
            \node[branch] (b1) [above right=0.2cm and 0.8cm of in] {RLS-ARX (Linéaire)\\{\tiny Poids légers (< 1 Ko RAM)}};
            \node[branch] (b2) [below right=0.2cm and 0.8cm of in] {GRU (Non-linéaire)\\{\tiny INT8 Quantifié (< 2 Mo RAM)}};
            \node[merge] (m) [below right=0.2cm and 0.8cm of b1] {Fusion};
            \node[output] (out) [right=0.6cm of m] {Prédiction};

            \draw[-latex, thick] (in) |- (b1);
            \draw[-latex, thick] (in) |- (b2);
            \draw[-latex, thick] (b1) -| (m);
            \draw[-latex, thick] (b2) -| (m);
            \draw[-latex, thick] (m) -- (out);
        \end{tikzpicture}
      \end{exampleblock}
    \end{column}
  \end{columns}
\end{frame}

% Slide 5: Le Workflow d'Entraînement pas à pas
\begin{frame}{3. Le Workflow Technique d'Entraînement}
  Le passage du modèle de la phase d'apprentissage globale à l'exécution en ligne locale suit 5 étapes clés :
  
  \vspace{0.8em}
  \centering
  \begin{tikzpicture}[scale=0.75, every node/.style={transform shape},
      stepbox/.style={rectangle, draw=blue!50!black, fill=blue!10, rounded corners, minimum width=2.4cm, minimum height=1.1cm, align=center, font=\scriptsize, text width=2.2cm},
      arrow/.style={-latex, thick, blue!50!black}]
      
      \node[stepbox] (s1) {1. Pré-entraînement\\{\tiny Global (Float32)\\[0.2em] sur SETI@home}};
      \node[stepbox] (s2) [right=0.4cm of s1] {2. Quantification\\{\tiny TinyML (INT8)\\[0.2em] Division RAM par 4}};
      \node[stepbox] (s3) [right=0.4cm of s2] {3. Déploiement\\{\tiny Agent Local sur PC volontaires}};
      \node[stepbox] (s4) [right=0.4cm of s3] {4. Adaptation\\{\tiny RLS-ARX récursive\\[0.2em] minute par minute}};
      \node[stepbox] (s5) [right=0.4cm of s4] {5. Fine-tuning\\{\tiny GRU-EWC sur 72h\\[0.2em] toutes les 24h}};
      
      \draw[arrow] (s1) -- (s2);
      \draw[arrow] (s2) -- (s3);
      \draw[arrow] (s3) -- (s4);
      \draw[arrow] (s4) -- (s5);
  \end{tikzpicture}
\end{frame}

% Slide 6: Étape 1 : Pré-entraînement sur SETI@home
\begin{frame}{4. Étape 1 : Pré-entraînement sur le dataset global}
  \begin{columns}[T]
    \begin{column}{0.48\textwidth}
      \begin{block}{Pourquoi cette étape ?}
        \begin{itemize}
          \item Les réseaux de neurones profonds ont besoin de beaucoup de données pour converger.
          \item Les traces locales de Yaoundé sont encore trop jeunes ou limitées.
          \item Nous utilisons la masse de données de **SETI@home** pour apprendre au GRU les comportements de base des ordinateurs.
        \end{itemize}
      \end{block}
    \end{column}
    
    \begin{column}{0.48\textwidth}
      \begin{block}{Ce que le modèle apprend}
        \begin{itemize}
          \item Les dynamiques d'activité utilisateur (heures de repas, repos nocturne).
          \item Les cycles d'extinction volontaire de fin de journée.
          \item Les relations temporelles d'une heure à l'autre.
        \end{itemize}
      \end{block}
    \end{column}
  \end{columns}
\end{frame}

% Slide 7: Étape 2 & 3 : Quantification TinyML et Déploiement
\begin{frame}{5. Étape 2 \& 3 : Quantification et Déploiement Frugal}
  \begin{columns}[T]
    \begin{column}{0.48\textwidth}
      \begin{block}{La Quantification INT8}
        \begin{itemize}
          \item Conversion des poids mathématiques de Float32 (4 octets) à INT8 (1 octet).
          \item \textbf{Impact direct :} La RAM requise pour le GRU passe de **8 Mo à moins de 2 Mo**.
          \item Préservation du score de prédiction : le modèle conserve une précision de $F1 \ge 0,80$.
        \end{itemize}
      \end{block}
    \end{column}
    
    \begin{column}{0.48\textwidth}
      \begin{block}{Le Déploiement Frugal}
        \begin{itemize}
          \item L'agent local n'embarque pas de gros frameworks.
          \item Utilisation d'un moteur d'inférence minimaliste (ex: TensorFlow Lite Micro).
          \item Zéro perturbation pour l'utilisateur volontaire (utilisation CPU inférieure à 0,1\%).
        \end{itemize}
      \end{block}
    \end{column}
  \end{columns}
\end{frame}

% Slide 8: Étape 4 & 5 : Adaptation Récursive et Fine-Tuning EWC
\begin{frame}{6. Étape 4 \& 5 : Adaptation Locale et Anti-Oubli}
  Une fois sur la machine de l'étudiant à Yaoundé, le modèle doit s'adapter aux pannes d'électricité locales :
  
  \vspace{0.5em}
  \begin{columns}[T]
    \begin{column}{0.48\textwidth}
      \begin{block}{Adaptation Instantanée (RLS-ARX)}
        Le modèle linéaire met à jour ses coefficients toutes les minutes dès qu'une coupure survient, en associant la baisse de tension ou le passage sur batterie à la perte imminente de disponibilité.
      \end{block}
    \end{column}
    
    \begin{column}{0.48\textwidth}
      \begin{block}{Fine-Tuning Local Périodique (GRU)}
        \begin{itemize}
          \item Le GRU est ré-entraîné toutes les 24 heures sur le buffer local.
          \item **Régularisation EWC :** Bloque la modification des poids critiques appris sur SETI@home pour éviter l'oubli catastrophique des cycles généraux.
        \end{itemize}
      \end{block}
    \end{column}
  \end{columns}
\end{frame}

% Slide 9: Intégration dans l'Ordonnanceur (Scheduler)
\begin{frame}{7. Comment la prédiction est-elle exploitée ?}
  La prédiction de disponibilité alimente directement le middleware de calcul volontaire :
  
  \vspace{0.5em}
  \begin{columns}[T]
    \begin{column}{0.48\textwidth}
      \begin{block}{Soumission Intelligente des Tâches}
        \begin{itemize}
          \item Avant d'envoyer un calcul à un nœud volontaire, l'ordonnanceur central interroge l'agent local.
          \item Si la disponibilité estimée sur la durée du calcul est élevée ($\ge 0,80$), la tâche est soumise.
          \item Permet d'éviter d'envoyer des tâches longues sur des nœuds sur le point de subir un délestage.
        \end{itemize}
      \end{block}
    \end{column}
    
    \begin{column}{0.48\textwidth}
      \begin{block}{Migration Préventive}
        \begin{itemize}
          \item Si le modèle local détecte un délestage imminent (ex: baisse continue de la batterie), il alerte l'ordonnanceur.
          \item Le calcul en cours est sauvegardé (checkpoint) et migré vers une machine plus stable du réseau.
        \end{itemize}
      \end{block}
    \end{column}
  \end{columns}
\end{frame}

% Slide 10: Conclusion et Perspectives
\begin{frame}{8. Conclusion}
  \begin{block}{Une Approche Rigoureuse et Frugale}
    L'architecture hybride et le workflow d'apprentissage de VC-UY1 apportent des réponses concrètes aux défis locaux :
    \begin{itemize}
      \item \textbf{Modélisation Hybride :} Le RLS-ARX pour la réactivité électrique, le GRU-EWC pour les habitudes humaines.
      \item \textbf{Workflow Scientifique :} Pré-entraînement mondial (SETI@home) garantissant une base stable, suivi d'une adaptation locale sécurisée.
      \item \textbf{Frugalité Validée :} Utilisation de la quantification INT8 TinyML permettant de respecter les limites strictes du matériel local ($< 10$ Mo RAM).
    \end{itemize}
  \end{block}
  
  \centering
  \vspace{1em}
  \Large\textbf{Merci pour votre attention. Place aux questions.}
\end{frame}

\end{document}
